% ============================================================
% 00-abstract.tex
% ============================================================
\begin{abstract}
Student counseling in higher education draws on heterogeneous institutional knowledge, including relational curricula, cohort-specific tuition schedules, statutory exemption criteria, scholarship policies, and narrative administrative regulations. This setting poses a challenge for uniform retrieval and reasoning pipelines, as different forms of institutional knowledge may benefit from different representations, retrieval mechanisms, and processing strategies. We present CTU-Chat, a centralized supervisor-routed, domain-specialized multi-agent system for Vietnamese university counseling. The supervisor directs each query to one of four specialists, while each specialist is assigned an appropriate knowledge source, retrieval path, and set of tools. Relational academic information is handled through a knowledge graph; structured tuition and scholarship information is supported by dedicated lookup and calculation functions; and narrative regulations are retrieved through hybrid document search. Evaluation shows 97.0\% raw structured-supervisor routing accuracy. On a redesigned 100-query held-out test set, the full retrieval stack achieves 0.7400 Hit@1, 0.8608 Recall@5, and 0.8333 MRR@10, while the full end-to-end generation reaches 0.6534 Context Recall and 0.8636 Faithfulness with 0.8325 Source Recall. These findings suggest that centralized coordination with domain-specific knowledge allocation and multi-representation evidence packing can support heterogeneous university-counseling tasks under the evaluated conditions.

\keywords{Multi-Agent Systems \and Domain-Specialized Agents \and Centralized Supervisor Routing \and Heterogeneous Knowledge Allocation \and Retrieval-Augmented Generation \and University Counseling}
\end{abstract}
